\section{Method}
\label{sec:method}

\subsection{Problem setup}

A PuzzleScript game is a tuple $(\mathcal{O}, R, \{L_i\})$ where
$\mathcal{O}$ is an ordered list of object kinds, $R$ is a list of
rewrite rules, and $\{L_i\}$ are authored levels. Each level is a
multi-channel grid $s \in \{0,1\}^{C \times H \times W}$ where channel
$c$ is the indicator of object kind $c$ at each cell. An action $a \in
\{\texttt{up}, \texttt{down}, \texttt{left}, \texttt{right},
\texttt{action}\}$ together with the current state $s_t$
deterministically yields a next state $s_{t+1}$ via the engine. The
state can also carry a terminal \emph{won} bit, since some rules
trigger a win on certain conditions.

We train a model $f_\theta(s_t, a_t, \mathrm{spec}) \mapsto
\widehat{s}_{t+1}$ where $\mathrm{spec}$ is the tokenized game source
(a sequence of integers from a shared $\sim$140-symbol vocabulary that
covers the engine's grammar plus per-game object names and 5$\times$5
sprite RGBA values when sprite tokens are enabled). At evaluation we
roll the model out autoregressively and compare cell-by-cell against
the engine.

\subsection{Rule-conditioned NCA world model}

\paragraph{Slot encoder.}
A perceiver-style encoder $\mathrm{Enc}_\phi$ maps the game's token
sequence $\mathrm{spec} \in \mathbb{Z}^{S}$ to $K$ slot vectors $z \in
\mathbb{R}^{K \times d}$. The encoder applies $\ell$ self-attention
layers over the tokens, then $K$ learned query vectors cross-attend
into the contextualized token sequence to produce $z$. Slots are
shared across cells and remain fixed throughout the rollout.

\paragraph{NCA dynamics body.}
We embed the input $(s_t, a_t)$ into a hidden grid $h^{(0)} \in
\mathbb{R}^{H \times W \times d_h}$ via a single dense layer applied
to the per-cell concatenation of $s_t$'s channels and a broadcast
one-hot action. We then apply $T$ NCA steps:
\begin{align}
  h^{(k+1)} &= h^{(k)} + W_o \,\mathrm{GELU}\!\bigl(
      W_c\, [\, h^{(k)} \,;\, h^{(0)} \,]_{3\times 3}
      \;+\;
      \mathrm{XAttn}(h^{(k)}, z)
    \bigr) \, ,
\end{align}
where $[\,\cdot\,]_{3\times 3}$ denotes a $3{\times}3$ convolution
producing a vector at each cell, $\mathrm{XAttn}$ is a multi-head
attention with each cell as a query and the $K$ slots as keys/values,
and $W_o, W_c$ are shared across all $T$ steps. The
\emph{input-skip} concatenation re-injects the embedded
$(s_t, a_t)$ at every step, preventing the deep unroll from drifting.

\paragraph{Pool features.}
A bare $3{\times}3$ convolution can only propagate one cell of context
per NCA step, but several PuzzleScript rule patterns reference cells
that are arbitrarily far apart (\texttt{[ X | ... | Y ]} along an
axis, or \texttt{[X][Y]} anywhere in the grid). We therefore augment
each step with three optional pool feature stacks:
axis max pool (row-max and col-max), axis cumulative max
(directional prefix-max in each of four directions), and global max
pool. These features are computed from the current hidden grid and
broadcast back to every cell before the convolution, exposing
non-local information that would otherwise require $\Omega(\max(H,W))$
NCA steps to surface.

\paragraph{Hierarchical weight sharing.}
Inspired by the engine's two natural levels of iteration (an inner
ordered list of rules + an outer \texttt{again} re-evaluation loop),
we factor $T = n_\ell \cdot n_r$ where $n_\ell$ distinct layers form
the inner block and $n_r$ outer applications share weights. Setting
$n_r = T$ recovers a single shared layer applied $T$ times (matches
the engine's outer-loop structure); setting $n_r = 1$ recovers the
usual per-step body with $T$ distinct layers.

\paragraph{Heads.}
A shared \texttt{readout} dense layer maps the final $h^{(T)}$
cell-wise to next-state logits $\widehat{s}_{t+1} \in \mathbb{R}^{C
\times H \times W}$; a separate dense layer over the spatially
average-pooled $h^{(T)}$ produces a binary win logit. Training loss
is a per-cell BCE on the next-state with an upweight on cells that
\emph{change} between $s_t$ and $s_{t+1}$, plus a binary cross-entropy
on the win head.

\subsection{Baselines}

We compare against three non-NCA baselines that share the slot
encoder $\mathrm{Enc}_\phi$, the input embedding, and the readout
heads, but replace the NCA body. Each baseline consumes
$h^{(0)}$ and the slots $z$ and emits a final hidden grid $h^{(T)}$
of the same shape; the comparison isolates the contribution of the
spatial-update body.

\begin{itemize}
  \item \textbf{CNN}: a stack of $n_b$ residual blocks, each
    containing the same conv + optional pool + slot-cross-attn +
    residual gate as one NCA layer, but with \emph{distinct}
    weights per block and applied once. Tests whether iteration with
    weight sharing matters at all.
  \item \textbf{U-Net}: a 2-level encoder--decoder with skip
    connections and a FiLM bottleneck that reads pooled slots. Tests
    whether multi-scale hierarchy substitutes for repeated local
    updates.
  \item \textbf{ViT}: $n_v$ Transformer encoder layers over the
    flattened cell sequence with learned 2D positional embedding,
    plus a slot cross-attention and feed-forward block per layer.
    Tests whether unrestricted global attention beats locality plus
    iteration.
\end{itemize}

All baselines accept the same input $(s_t, a_t, \mathrm{spec})$ and
produce the same $(\widehat{s}_{t+1}, \widehat{w}_{t+1})$ outputs as
the NCA, allowing identical training and evaluation pipelines.

\subsection{Adaptive halting (optional)}
\label{sec:method-halt}

A fixed $T$ is wasteful for easy transitions and insufficient for
hard ones. We add an optional halting head: at each NCA step $k$, a
small dense layer reads the spatially pooled $h^{(k)}$ and emits a
halt logit $\lambda_k$. The induced halt distribution
$p_k = \mathrm{sigmoid}(\lambda_k) \prod_{j<k}(1 - \mathrm{sigmoid}(\lambda_j))$
weights the per-step loss
$\mathcal{L}_{\text{rec}} = \sum_k p_k \mathcal{L}_k$
plus a KL-to-geometric-prior regularizer
$\mathrm{KL}(p \,\|\, \mathrm{Geom}(p_0))$. We additionally study a
\emph{convergence-stopping} baseline that halts at inference when
consecutive predictions stop changing, with no learned halt head.
